% Exercise 4.1-5
\begin{algorithmic}[1]
  \TITLE{\textsc{Find-Maximum-Subarray}$(A, \mathit{low}, \mathit{high})$}
  \STATE $\mathit{maxLow} = \mathit{low}$
  \STATE $\mathit{maxHigh} = \mathit{low}$
  \STATE $\mathit{maxSum} = A[\mathit{low}]$
  \STATE $\mathit{currentLow} = \mathit{low}$
  \STATE $\mathit{currentSum} = A[\mathit{low}]$
  \FOR{$j = \mathit{low} + 1$ \TO $\mathit{high}$}
  \IF{$\mathit{currentSum} > 0$}
  \STATE $\mathit{currentSum} = \mathit{currentSum} + A[j]$
  \ELSE
  \STATE $\mathit{currentLow} = j$
  \STATE $\mathit{currentSum} = A[j]$
  \ENDIF

  \IF{$\mathit{currentSum} > \mathit{maxSum}$}
  \STATE $\mathit{maxSum} = \mathit{currentSum}$
  \STATE $\mathit{maxLow} = \mathit{currentLow}$
  \STATE $\mathit{maxHigh} = j$
  \ENDIF
  \ENDFOR

  \RETURN $(\mathit{maxLow}, \mathit{maxHigh}, \mathit{maxSum})$
\end{algorithmic}

The \textbf{for} loop of lines 6--16 extends the maximum subarray ending at
$j - 1$ by $A[j]$ whenever that running sum is still positive, and otherwise
begins a fresh subarray at $j$. Since each element is examined exactly once,
the running time is $\Theta(n)$.

Note that $\mathit{maxSum}$ is initialised to $A[\mathit{low}]$ rather than to
the sum of the first two elements, so that a single element is a valid answer
and an array of entirely negative values returns its largest element.
